%==============================================================================
%  ABSTRACT / RÉSUMÉ  (front matter, unnumbered)
%==============================================================================
\chapter*{Abstract}
\addcontentsline{toc}{chapter}{Abstract}
\markboth{Abstract}{Abstract}

Machine learning models rarely fail because of their mathematics; they fail
because of everything around them, environments that cannot be reproduced,
experiments that are lost, and deployments that take weeks. This end-of-studies
project, carried out within \textbf{INSOMEA}, addresses this gap by designing and
implementing a complete, self-hosted \textbf{MLOps platform} that automates the
machine-learning lifecycle from a single web interface. The platform lets a data
scientist authenticate, organize work into projects, upload code and data, run
training on a Kubernetes cluster with automatic experiment tracking and artifact
browsing, register and promote model versions, deploy them as live prediction
services with automatic framework detection, expose them to external applications
through secured public API keys, and follow their behavior through per-prediction
serving telemetry (volume, latency, error rate), while an administrator oversees
users and resources. The system is built with FastAPI, Angular, Kubeflow
Pipelines, MLflow, KServe, MinIO, and PostgreSQL, is fully containerized, and is
continuously deployed to Azure Kubernetes Service. The work was managed with the
Scrum method over four sprints spanning six months, and validated end to end
with a real telecom churn use case.

\vspace{0.6cm}
\noindent\textbf{Keywords:} MLOps, Kubernetes, machine learning, experiment
tracking, model deployment, MLflow, KServe, FastAPI, Angular, Scrum.

\vspace{1.2cm}
\begin{otherlanguage}{french}
\section*{R\'esum\'e}
\addcontentsline{toc}{chapter}{R\'esum\'e}

Les mod\`eles d'apprentissage automatique \'echouent rarement \`a cause de leurs
math\'ematiques ; ils \'echouent \`a cause de tout ce qui les entoure, des
environnements impossibles \`a reproduire, des exp\'eriences perdues et des
d\'eploiements qui prennent des semaines. Ce projet de fin d'\'etudes, r\'ealis\'e
au sein d'\textbf{INSOMEA}, r\'epond \`a ce probl\`eme par la conception et la
r\'ealisation d'une \textbf{plateforme MLOps} compl\`ete et auto-h\'eberg\'ee qui
automatise le cycle de vie de l'apprentissage automatique depuis une interface web
unique. La plateforme permet \`a un data scientist de s'authentifier, d'organiser
son travail en projets, de t\'el\'everser code et donn\'ees, de lancer des
entra\^inements sur un cluster Kubernetes avec un suivi automatique des
exp\'eriences et la consultation des artefacts, d'enregistrer et de promouvoir des
versions de mod\`eles, de les d\'eployer comme services de pr\'ediction avec
d\'etection automatique du framework, de les exposer \`a des applications externes
via des cl\'es API publiques s\'ecuris\'ees, et de suivre leur comportement par une
t\'el\'em\'etrie de service par pr\'ediction (volume, latence, taux d'erreur),
tandis qu'un administrateur supervise les utilisateurs et les ressources. Le syst\`eme repose sur FastAPI,
Angular, Kubeflow Pipelines, MLflow, KServe, MinIO et PostgreSQL, est enti\`erement
conteneuris\'e et d\'eploy\'e en continu sur Azure Kubernetes Service. Le travail a
\'et\'e conduit avec la m\'ethode Scrum en quatre sprints sur six mois, et
valid\'e de bout en bout par un cas d'usage r\'eel de churn t\'el\'ecom.

\vspace{0.6cm}
\noindent\textbf{Mots cl\'es :} MLOps, Kubernetes, apprentissage automatique,
suivi d'exp\'eriences, d\'eploiement de mod\`eles, MLflow, KServe, FastAPI,
Angular, Scrum.
\end{otherlanguage}
